% Part: lambda-calculus
% Chapter: introduction
% Section: fixed-point-combinator

\documentclass[../../../include/open-logic-section]{subfiles}

\begin{document}

\olfileid[pt]{lam}{int}{fix}
\olsection{Combinadores de ponto fixo}

Suponha que você tenha um termo lambda $g$ e queira outro termo $k$ com a
propriedade de que $k$ é $\beta$-equivalente a $gk$. Defina os termos
\[
\fn{diag}(x) = xx
\]
e
\[
l(x) = g(\fn{diag}(x))
\]
usando nossas convenções de notação; em outras palavras, $l$ é o termo
$\lambd[x][g(xx)]$. Seja $k$ o termo $ll$. Então temos
\begin{align*}
k & = (\lambd[x][g(xx)])(\lambd[x][g(xx)]) \\
& \red  g((\lambd[x][g(xx)])(\lambd[x][g(xx)])) \\
& = gk.
\end{align*}
Se tomarmos
\[
Y = \lambd[g][((\lambd[x][g(xx)])(\lambd[x][g(xx)]))]
\]
então $Yg$ e $g(Yg)$ se reduzem a um termo comum; logo $Yg \equiv_\beta
g(Yg)$. Isso é conhecido como ``combinador de Curry''. Se, em vez disso,
tomarmos
\[
Y = (\lambd[xg][g(xxg)])(\lambd[xg][g(xxg)])
\]
então, de fato, $Yg$ se reduz a $g(Yg)$, que é uma afirmação mais forte.
Essa última versão de $Y$ é conhecida como ``combinador de Turing''.

\end{document}

